The \peripheral{NPU} is a fixed-point neural network processing unit built around a single multiply-accumulate datapath and a configurable sequencer. One engine provides \textbf{four datapath modes}, selected per computation by the \texttt{NPUMODE} field (\register{NPUCR} bits 22:20): a fully-connected (dense) MLP layer (mode 0, the legacy datapath and the reset default), a one-dimensional convolution (mode 1), a binary XNOR-popcount layer (mode 2), and a general matrix--matrix multiply (mode 3). The peripheral operates on data staged in the shared NPU staging RAM (the 16\,KiB SRAM occupying \texttt{0xC000}--\texttt{0xFFFF}), which is multiplexed between the harts (reaching it through the multi-core arbiter) and the NPU's own compute port. Software is responsible for staging operands into the staging RAM, configuring the NPU, and reading the results after the computation is complete.

The NPU's \textit{register bus} is in the shared peripheral page (see Section \ref{s:multicore}), so any hart may configure the NPU, start a computation, and poll for completion. The \textit{data path} is likewise shared: all operands live in the shared NPU staging RAM at \texttt{0xC000}--\texttt{0xFFFF}, and any hart may stage operands there or retrieve results through the arbiter. While a computation is running, the in-NPU port multiplexer switches the staging RAM over to the NPU's compute port for the duration of the \texttt{THINK}, and \textbf{no hart is put to sleep}. Because the staging RAM belongs to the NPU during a run, software of any hart \textbf{must not read or write \texttt{0xC000}--\texttt{0xFFFF} while a computation may be active}: poll \register{NPUCR} bit 16 (\texttt{NPUTHINK}) and wait for it to read 0 before touching the staging RAM.

\subsubsection*{Data Formats}

All non-binary NPU operands are fixed-point numbers, each stored in its own 32-bit SRAM word:

\begin{itemize}
    \item \textbf{Inputs} (MLP/convolution inputs, GEMM A elements): signed Q0.24 format (25 bits: 1 sign bit + 24 fractional bits), representing values in $[-1, 1)$. Stored in bits 24:0 of a 32-bit SRAM word; bits 31:25 are ignored on read. Bit 24 is the sign bit.
    \item \textbf{Weights} (MLP/convolution weights, GEMM B elements): signed Q7.24 format (32 bits: 1 sign bit + 7 integer bits + 24 fractional bits), representing values in $[-128, 128)$.
    \item \textbf{Outputs}: signed Q7.24 format by default; the activation functions produce narrower ranges (see \textit{Activation Functions} below). XNOR mode packs its 1-bit operands 32 per word (see mode 2 below) and emits $\pm 1.0$ Q7.24 output words.
\end{itemize}

The accumulator is Q7.24 with \textbf{per-step round-to-nearest and saturation at $\pm 128$}: every multiply-accumulate step rounds and clips individually, in the fixed sequencer order. Sums whose intermediate values exceed $\pm 128$ clip; this is benign for normalized data but is a real ceiling for long dot products with large coefficients.

\subsubsection*{Datapath Modes}

The count registers \texttt{NPUNI} and \texttt{NPUNN} (\register{NPUCR} bits 15:8 and 7:0, each holding count~$-$~1) and the per-mode configuration words \register{NPUCFG1}/\register{NPUCFG2} are reinterpreted by each mode; the start-address registers \texttt{NPUIVSAR}/\texttt{NPUWVSAR}/\texttt{NPUOVSAR} always hold \textit{word indices within the staging RAM} (byte offset from \texttt{0xC000} divided by 4).

\begin{itemize}
    \item \textbf{Mode 0: MLP (legacy, reset default).} One dense layer: \texttt{NPUNI}~+~1 inputs at \texttt{NPUIVSAR}, \texttt{NPUNN}~+~1 output neurons. The weight matrix at \texttt{NPUWVSAR} holds one contiguous row per neuron; if bias is enabled (\texttt{NPUBEN}~=~1) each row begins with one bias weight (multiplied by an implicit input of 1.0) followed by the \texttt{NPUNI}~+~1 synaptic weights. \register{NPUCFG1}/\register{NPUCFG2} are unused.
    \item \textbf{Mode 1: 1-D convolution.} \texttt{NPUNI}~=~taps~$K - 1$, \texttt{NPUNN}~=~filters~$C_{out} - 1$. \register{NPUCFG1} packs stride, dilation, input length, and input channel count; \register{NPUCFG2} holds the host-computed output length $L_{out}$ (the sequencer \textit{trusts} it; the `valid'/`same' padding convention is a firmware decision). Inputs are channel-major at \texttt{NPUIVSAR}; weights are filter-major (bias first per filter when \texttt{NPUBEN}~=~1); outputs are written filter-major. See the register descriptions for the exact field packing.
    \item \textbf{Mode 2: XNOR-popcount (binary).} 1-bit activations and weights packed 32 per word, LSB-first (bit $i$ of the layer lives at bit $i \bmod 32$ of word $i \div 32$). \texttt{NPUNI}~=~packed words per neuron~$-$~1; the \textit{exact} bit count $K \le 4096$ lives in \register{NPUCFG2} bits 12:0, and hardware masks the unused tail bits of each last packed word out of the popcount. A neuron fires when $2 \cdot \mathrm{popcount}(\mathrm{XNOR}(a, w)) - K \ge$ the signed threshold in \register{NPUCFG1}, emitting $+1.0$ (\texttt{0x01000000}) else $-1.0$ (\texttt{0xFF000000}). \texttt{NPUBEN} and \texttt{NPUAEN} must be 0 in this mode.
    \item \textbf{Mode 3: GEMM.} $C = A \times B$ with $A$ ($M \times K$, row-major at \texttt{NPUIVSAR}, Q0.24 elements in $[-1,1)$), $B$ ($K \times N$, \textbf{column-major} at \texttt{NPUWVSAR}, Q7.24), and $C$ ($M \times N$, row-major at \texttt{NPUOVSAR}). $K-1$ in \texttt{NPUNI}, $N-1$ in \texttt{NPUNN}, $M-1$ in \register{NPUCFG1} bits 7:0. The working set $M{\cdot}K + K{\cdot}N + M{\cdot}N$ must fit the 4096-word staging RAM; larger problems tile by re-pointing the start-address registers between computations (there is \textbf{no hardware accumulation across tiles}, so $K$ must fit in one run). Because $C$ is row-major, a result matrix chains directly as the next layer's $A$ operand with no repacking (activated outputs are valid Q0.24 inputs; see below). \texttt{NPUBEN} must be 0 in this mode.
\end{itemize}

\textbf{The sequencer trusts its configuration}: there are no hardware bounds checks. A working set that exceeds the 4096-word staging RAM, or count/length fields inconsistent with the staged data, wraps the 12-bit word addresses and silently corrupts neighbouring staging-RAM regions; the engine never hangs, but the results (and possibly unrelated staged data) are garbage. Firmware owns operand sizing and tiling.

\subsubsection*{Activation Functions}

When \texttt{NPUAEN}~=~1, each output is passed through the activation selected by \texttt{NPUACTF} (\register{NPUCR} bits 25:23) before being written; \texttt{NPUAEN}~=~0 writes the raw Q7.24 accumulator regardless of \texttt{NPUACTF}. All activations share the one hardware sigmoid approximator (a piecewise-quadratic logistic function):

\begin{itemize}
    \item \textbf{0, sigmoid} (reset default): $\sigma(x)$, output Q0.24 in $[0,1)$; saturates for $|x| \ge 4$.
    \item \textbf{1, ReLU}: $\max(0, x)$, output Q7.24. Note that a ReLU output at or above 1.0 is \textit{not} a valid Q0.24 next-layer input; chain through normalized weights or use the clamp.
    \item \textbf{2, tanh} approximation, computed as $2\sigma(2x) - 1$ through the same sigmoid hardware: output in $[-1,1)$, sign-extended, whose low 25 bits are a valid Q0.24 next-layer input. Saturates for $|x| \ge 2$ (half the sigmoid's linear region, a consequence of the input doubling).
    \item \textbf{3, clamp} (hardtanh): the identity saturated to the Q0.24 rails $[-1, 1)$, sign-extended, a bounded linear activation whose output always chains as a Q0.24 input.
    \item \textbf{4, exponential approximation} $\widetilde{\exp}(x) = 2\sigma(x)$, output Q1.24 in $[0,2)$, for the softmax leg: hardware emits the per-element scores only, and \textbf{firmware performs the sum and normalization}. Subtracting the maximum logit before the run is recommended; the maximum element then maps to exactly 1.0.
    \item Codes 5--7 are reserved and behave as 0 (sigmoid).
\end{itemize}

The activation selection is latched when \texttt{THINK} is set, so a mid-computation write to \texttt{NPUACTF} (or \texttt{NPUMODE}) takes effect at the \textit{next} computation. XNOR mode (mode 2) ignores the activation path entirely; its $\pm 1.0$ encoding is fixed.

\subsubsection*{Operation}

\begin{enumerate}
    \item Stage the operands into the staging RAM per the selected mode's layout and write their word indices to \texttt{NPUIVSAR} and \texttt{NPUWVSAR}.
    \item Write the desired output word index to \texttt{NPUOVSAR}.
    \item Write the per-mode configuration (\register{NPUCFG1}/\register{NPUCFG2}) if the mode uses it.
    \item Configure \register{NPUCR}: \texttt{NPUNI}, \texttt{NPUNN}, \texttt{NPUMODE}, \texttt{NPUBEN}, \texttt{NPUAEN}, and \texttt{NPUACTF} as required by the mode.
    \item Set \texttt{NPUTHINK} (\register{NPUCR} bit 16) to~1 to start the computation. The staging RAM is now owned by the NPU; no hart may access \texttt{0xC000}--\texttt{0xFFFF} until the computation completes.
    \item Poll \register{NPUCR} bit 16 (\texttt{NPUTHINK}) until it reads~0, indicating that the computation is complete and the staging RAM has been released back to the harts, or take the think-done interrupt (below).
    \item Read the outputs from the staging RAM at the word index stored in \texttt{NPUOVSAR}.
\end{enumerate}

Run times scale with the mode's loop product (e.g.\ a GEMM costs on the order of $5 \cdot M \cdot N \cdot K$ NPU clock cycles; a convolution $5 \cdot C_{out} \cdot L_{out} \cdot C_{in} \cdot K$). Full-scope convolutions can reach hundreds of milliseconds at 24\,MHz; size polling budgets (or use the think-done interrupt) accordingly. XNOR mode is the cheapest per output by roughly an order of magnitude (32 synapses per fetched word pair).

\subsubsection*{Think-done interrupt (vector 120)}

As an alternative to polling, the NPU raises the \textit{think-done} interrupt (interrupt vector 120) when a computation completes. The \texttt{NPUTHINKDONE} flag (\register{NPUSR} bit 0) is set once per computation, by the same internal event that self-clears \texttt{NPUTHINK}, and is \textit{sticky}: it remains set until software clears it by writing a 1 to \register{NPUSR} bit 0 (write-1-to-clear; writing 0 has no effect, and reading \register{NPUSR} never clears it). The interrupt is a level, asserted while both \texttt{NPUTHINKDONE} and the enable bit \texttt{NPUTDIE} (\register{NPUCR} bit 19) are set; it is routed, claimed, and completed through the interrupt router like every other peripheral source. \texttt{NPUTDIE} resets to 0, so software that polls is unaffected by the interrupt's existence. Every mode uses the same completion event and the same vector.

A typical interrupt-driven flow arms \texttt{NPUTDIE}, starts the computation, and sleeps; the handler clears \texttt{NPUTHINKDONE} (write-1-to-clear) before completing the interrupt at the router. \textbf{The staging-RAM ownership contract is unchanged by the interrupt}: the staging RAM belongs to the NPU while a computation may be active, and no hart may access \texttt{0xC000}--\texttt{0xFFFF} until \register{NPUCR} bit 16 (\texttt{NPUTHINK}) reads 0. Taking the think-done interrupt implies the computation has finished, but any hart \textit{other} than the one servicing the interrupt must still consult \texttt{NPUTHINK} before touching the staging RAM.
